\documentstyle []{article}
\oddsidemargin -7mm
\evensidemargin -7mm
\topmargin -05mm
\headheight 0pt
\headsep 0pt
\textheight 235mm
\textwidth 174mm
\pagestyle{empty}

\begin{document}

\title{Issues in Symbolic Parametrization
       of Algebraic Curves 
\footnote{Supported by 
{\it \"Osterr. Fonds zur F\"orderung der wissenschaftlichen
Forschung}, Proj. POSSO, Nr. P9181-TEC. (ESPRIT BRA No. 6846).}  }

\author{Franz  Winkler \\
        Institut f\"ur Mathematik and RISC-LINZ \\
        Johannes Kepler Universit\"at \\
        A-4040 Linz, Austria \\
        e-mail: {\tt winkler@risc.uni-linz.ac.at}}

\date{ }

\maketitle

\thispagestyle{empty}

\vspace{4 mm}

\noindent{\bf Abstract}

\vspace{5 mm}

The problem of rational parametrization of plane algebraic curves is best
explained by an example. 
Suppose we consider the curve ${\cal C}_1$ in the complex plane, whose
points are the zeros of the polynomial
$$f_1(x,y) = (x^2+4y+y^2)^2 - 16(x^2+y^2).$$
${\cal C}_1$ is a curve of genus 0, so it can be rationally 
parametrized, e.g. as
\begin{eqnarray*}
  & x(t) =
          { 9216t^4-62464t^3+141696t^2-108864t+2916 \over
           430336t^4-94464t^3+58976t^2-5904t+1681        }\ , \\
  & y(t) =
          { 40960t^4-184320t^3+204800t^2+11520t-12960 \over
           430336t^4-94464t^3+58976t^2-5904t+1681        }\ .
\end{eqnarray*}


\vspace{2 mm}

In general we say that
a curve $\cal C$ has a {\it rational parametrization}, iff there are 
(nontrivial)
rational functions $x(t),y(t)$ such that the defining polynomial $f$ of
$\cal C$ vanishes on them,
$$f(x(t),y(t))=0.               $$ 

\vspace{2 mm}

The parametrizable curves are exactly the irreducible curves of
genus 0, i.e. they achieve the upper bound on the number of 
singularities.
If such a parametrization exists, then it
is by no means unique, even if we require that the degrees of the
rational functions involved are minimal. In fact, one of the main 
problems is to determine a minimal extension of the field of definition,
in which a parametrization can be expressed.

For this purpose we assume that $K$ is a computable field of
characteristic 0, we call it the {\it field of definition}.
$\overline K$ denotes the algebraic closure of $K$.
The defining polynomial of the curve $\cal C$ that we want to
parametrize has coefficients in $K$.
The curve $\cal C$ itself, however, is considered to be a curve
over $\overline K$.
In the parametrization process  we might have to
extend $K$ algebraically. We will ultimately construct a
parametrization (if one exists) over some field 
$K(\gamma) \subseteq \overline K$, where $\gamma$ is algebraic
over $K$.
Of course, we want to keep  the degree of $\gamma$ as low as possible.

\vspace{2 mm} 

Singular points on algebraic curves occur only in full conjugacy  classes.
Consequently they do not introduce algebraic numbers in the
parametrization process. In general, however, we also need 
a certain number of simple points on the curve. This might force us
to extend the field of definition of the curve, in order to
arrive a field over which we can give a parametrization.


From the work of N\"other              and
Hilbert and Hurwitz                 we
know that it is possible to parametrize any curve $\cal C$ of genus 0
over the field of definition $K$, if $\deg({\cal C})$ is odd,
and over some quadratic extension of $K$, if $\deg({\cal C})$ is even.
In a previous paper J.R. Sendra and the author have presented an 
algorithm which actually achieves this optimal field of 
parametrization.
Moreover, if the field of definition is the field of rational
numbers, we can also
decide if the curve can be parametrized over the reals, and compute
such a parametrization if it exists.

\vspace {2 mm}

Our parametrization algorithm is of polynomial bit complexity
if we assume that the curve $\cal C$ defined over the rational numbers
has only ordinary singularities.
If we allow curves with also non--ordinary singularities, then
the bit complexity 
becomes exponential in the worst case.
\vspace{2 mm}
\newline
{\bf Keywords:} parametrization of algebraic curves, rational points on 
algebraic curves, complexity.




\end{document}
